\telubab{1}{LAMPIRAN A}{FULL AGENT PROMPT}{sec:A}

This appendix gives the complete prompt the agent receives, for one real
holdout question. The text was produced by
\texttt{build\_\allowbreak{}prompt}
(\texttt{src/\allowbreak{}needle/\allowbreak{}inference/\allowbreak{}runner.\allowbreak{}py}),
the same function used in training and evaluation. It is passed to
Qwen2.5-7B-Instruct as a single user message through the model's chat
template.

\subsection{How the example was
reproduced}\label{how-the-example-was-reproduced}

\begin{itemize}
\tightlist
\item
  \textbf{Question:}
  \texttt{5abd\allowbreak{}9452\allowbreak{}5542\allowbreak{}992a\allowbreak{}c4f3\allowbreak{}82d2\allowbreak{}}
  from the frozen holdout
  (\texttt{configs/\allowbreak{}evaluation/\allowbreak{}holdout\_\allowbreak{}v1\_\allowbreak{}ids.\allowbreak{}json}).
  The row was taken from the HotpotQA distractor validation split on
  Hugging Face (\texttt{hotpotqa/\allowbreak{}hotpot\_\allowbreak{}qa},
  \texttt{distractor/\allowbreak{}validation-\allowbreak{}00000-\allowbreak{}of-\allowbreak{}00001.\allowbreak{}parquet})
  and converted to the native HotpotQA format.
\item
  \textbf{History:} a \texttt{SearchEnvironment} (top-k = 3, budget = 3)
  replayed the two searches that R1 seed 8 made in its holdout episode
  (\texttt{results/\allowbreak{}2026-\allowbreak{}10-\allowbreak{}10-\allowbreak{}batch/\allowbreak{}r1\_\allowbreak{}full\_\allowbreak{}seed8/\allowbreak{}episodes.\allowbreak{}jsonl}).
  \texttt{build\_\allowbreak{}prompt} was then called at the point where
  the agent chose its final action.
\item
  \textbf{Gold supporting facts:} \texttt{2014 S/\allowbreak{}S|0} and
  \texttt{Winner (band)|0}.
\item
  \textbf{Agent's action at this point:}
  \texttt{ANSWER{[}\allowbreak{}YG Entertainment{]} CITATIONS{[}\allowbreak{}2014 S/\allowbreak{}S|0; Winner (band)|0{]}}.
  The answer is correct and both gold sentences are cited
  (Table~\ref{tbl:3.1}, episode A).
\end{itemize}

\subsection{Prompt text}\label{prompt-text}

The prompt is reproduced verbatim. The leading spaces on some sentences
are present in the HotpotQA source text.

\begin{Verbatim}
You are a careful evidence-seeking QA agent.
Return exactly one action: SEARCH[query] or ANSWER[answer] CITATIONS[title|index; ...].
SEARCH[query] is search-only: never append CITATIONS[...] to a SEARCH action.
ANSWER[...] is terminal; when the search budget reaches 0, answer immediately.
Example: ANSWER[United States] CITATIONS[Ursula K. Le Guin|1; The Left Hand of Darkness|1].
Use one title|index citation per entry; separate entries with semicolons, never comma-separated indices.
Do not add a period, explanation, or any text after the final closing bracket.
Copy citation titles exactly as shown in the search results, including punctuation and accents.
Never invent a document title; use CITATIONS[] when no citation can be validated.
Use only the supplied search results for citations.
Question: 2014 S/S is the debut album of a South Korean boy group that was formed by who?
Remaining search budget: 1
Search results so far:
Search 1: S/S 2014 debut South Korean boy group
[2014 S/S]
0: 2014 S/S is the debut album of South Korean group WINNER.
1:  It was released on August 12, 2014 by the group's record label, YG Entertainment.
2:  The members were credited for writing the lyrics and composing the majority of the album's songs.
[History (band)]
0: History (Korean: [Hangul] ) was a South Korean boy group formed by LOEN Entertainment in 2013.
1:  They debuted on April 26, 2013 with "Dreamer", featuring the narration of their labelmate IU.
2:  They were LOEN Entertainment's first boy group.
3:  They officially disbanded on May 12, 2017.
[SF9 (band)]
0: SF9 (Korean: [Hangul] ; shortened from Sensational Feeling 9) is a South Korean boy group formed by FNC Entertainment.
1:  SF9 is the company's first dance boy group to ever debut.
2:  SF9 debuted on October 5, 2016 with the release of their first single album "Feeling Sensation".
Search 2: WINNER YG Entertainment
[Winner (band)]
0: Winner (Hangul: [Hangul]), often stylized as WINNER, is a South Korean boy group formed in 2013 by YG Entertainment and debuted in 2014.
1:  It currently consists of four members, Jinwoo, Seunghoon, Mino and Seungyoon.
2:  Originally a five-piece group with Taehyun, who later departed from the group in November 2016.
[2014 S/S]
0: 2014 S/S is the debut album of South Korean group WINNER.
1:  It was released on August 12, 2014 by the group's record label, YG Entertainment.
2:  The members were credited for writing the lyrics and composing the majority of the album's songs.
[History (band)]
0: History (Korean: [Hangul] ) was a South Korean boy group formed by LOEN Entertainment in 2013.
1:  They debuted on April 26, 2013 with "Dreamer", featuring the narration of their labelmate IU.
2:  They were LOEN Entertainment's first boy group.
3:  They officially disbanded on May 12, 2017.
\end{Verbatim}

\subsection{Notes on the prompt}\label{notes-on-the-prompt}

\begin{itemize}
\tightlist
\item
  \textbf{First turn.} Before any search, the history section reads
  \texttt{(none)} and the remaining budget is 3.
\item
  \textbf{Repeated documents.} Each search lists its full top-3 results,
  even when a document was already returned by an earlier search (``2014
  S/S'' and ``History (band)'' appear twice above). The prompt grows
  with every search, which is one reason the budget is limited to 3.
\item
  \textbf{Why the format rules exist.} The rules about citation
  formatting (one \texttt{title|index} per entry, exact titles, no
  trailing text) target the protocol failures observed in the un-tuned
  model, which lost 37 of 300 holdout episodes to malformed actions
  (Section~\ref{sec:4.2}).
\item
  \textbf{No-citation loophole.} The line ``use CITATIONS{[}{]} when no
  citation can be validated'' is what makes an empty citation list a
  valid action (the loophole discussed in Section~\ref{sec:3.3}).
\end{itemize}
